\section{Постановка задачи}
\label{sec:task}

Целью лабораторной работы является ознакомление, анализ и приобретение навыков реализации модели линейной рециркуляционной сети для задачи сжатия графической (или иной) информации.

Вариант задания~\mdash  адаптивный коэффициент обучения и нормированные веса (\texttt{ADAPTIVE~=~True}, \texttt{NORMALIZE\_WEIGHTS~=~True}).

Требуется:
\begin{enumerate}
	\item построить модель линейной рециркуляционной сети, решающую задачу сжатия графических данных;
	\item построить семейства графиков зависимостей параметров модели;
	\item предоставить письменный отчёт.
\end{enumerate}

Исходные данные~\mdash  растровое изображение формата BMP. Изображение размера $h \times w$ разбивается на $Q$ прямоугольных блоков размера $r \times m$, где $4 \le r \le h$, $4 \le m \le w$. Значения интенсивности нормализуются в диапазон $[-1,\,1]$ по формуле
\begin{equation}
	c = \frac{2C}{C_{\max}} - 1,\qquad C_{\max} = 255.
\end{equation}
Каждый блок преобразуется в эталонный вектор $[X]$ длины $n = r \cdot m \cdot |S|$, где $|S|$~\mdash  число цветовых каналов.
